% =============================================================================
% 028 / 068
% Status: raw
% =============================================================================
% Notes:
%
% Source question:
% 如果語法並不能區別於官話的話，就算像香港、南粵一樣，很多字不用正確漢字用俗字、有粵語作為官方語言的香港，但只要香港的媒體和教育也推普，被殖民同化會非常快。而比如東南亞百越裔能熟練掌握多種語言，就是因為這些語言語法不同，代表的思維、場景不同。所以，如果吳語不能高階化，不能語法上跟官話隔絕，最終還是會被同化。
% 我的理解對嗎？
% =============================================================================

\chapter[如果語法並不能區別於官話的話，就算像香港、南粵一樣，很多字不用正確漢字用俗字、有粵語作為官方語言的香港，但只要香港的媒體和教育也推普，被殖民同化會非常快。而比如東南亞百越裔能熟練掌握多種語言，就是因為這些語言語法不同，代表的思維、場景不同。所以，如果吳語不能高階化，不能語法上跟官話隔絕，最終還是會被同化。 我的理解對嗎？]{如果語法並不能區別於官話的話，就算像香港、南粵一樣，很多字不用正確漢字用俗字、有粵語作為官方語言的香港，但只要香港的媒體和教育也推普，被殖民同化會非常快。而比如東南亞百越裔能熟練掌握多種語言，就是因為這些語言語法不同，代表的思維、場景不同。所以，如果吳語不能高階化，不能語法上跟官話隔絕，最終還是會被同化。 \\[0.35em] 我的理解對嗎？}
\qaanswer{
如果不能讓自己的母語言文一致，不能用書面語表達自己母語的語言思維。那麼就不會有標準出來。因為標準就是需要可以被量化、記錄、傳播的。任何一種沒有標準的規範化的語法和構詞法的口語是無法在強勢語言的滲透下存活的。單純透過情懷、社會呼籲和家庭堅守來保護吳語是蒼白和無力的。
}

